%%
%% part5-examples/ch29-physics-book.tex
%%
%% Chapter 29: A Physics Book — A complete sample chapter
%% from a physics textbook, focusing on equations, units,
%% and experimental data.
%%

\chapter{کتاب فیزیک}
\label{chap:physics-book}

\section{چرا این موضوع مهم است؟}
\label{sec:physics-book-why}

در دو فصل قبل، نمونه‌هایی از کتاب ریاضی
(با تمرکز بر قضیه و اثبات) و کتاب
برنامه‌نویسی (با تمرکز بر کد و الگوریتم)
دیدیم. در این فصل، نمونه‌ای از کتاب
فیزیک را می‌بینیم که چالش‌های خاص خود
را دارد: معادلات، یکاها، داده‌های تجربی
و نمودارها.

هدف این فصل، نشان دادن این است که چگونه
می‌توان از \lr{MatinBook} برای نوشتن
کتاب‌های فیزیک استفاده کرد.

\mnote{کتاب‌های فیزیک معمولاً ترکیبی از معادلات پیچیده، جداول داده‌های تجربی و نمودارهای دقیق هستند.}

\section{ساختار فصل نمونه}
\label{sec:physics-book-structure}

فصل نمونه، از یک کتاب فیزیک خیالی با
عنوان «مبانی مکانیک کلاسیک» است. این
فصل شامل بخش‌های زیر است:

\begin{itemize}
    \item مقدمه.
    \item معادلات حرکت.
    \item یکاها و ابعاد.
    \item داده‌های تجربی.
    \item نمودارها.
    \item تمرین‌ها.
\end{itemize}

\section{مقدمه}
\label{sec:physics-book-intro}

مکانیک کلاسیک\index{مکانیک کلاسیک}،
شاخه‌ای از فیزیک است که به مطالعه‌ی حرکت
اجسام می‌پردازد. این شاخه، پایه‌ی بسیاری
از شاخه‌های دیگر فیزیک است.

در این فصل، با قوانین نیوتن و معادلات
حرکت آشنا می‌شویم.

\mnote{مکانیک کلاسیک توسط \lr{Isaac Newton} در قرن هفدهم پایه‌گذاری شد.}

\section{قوانین نیوتن}
\label{sec:physics-book-newton}

\subsection{قانون اول}
\label{sec:physics-book-newton-first}

\begin{theorem}[قانون اول نیوتن]
    هر جسم در حال سکون یا حرکت یکنواخت
    در خط راست، به همان حالت باقی می‌ماند
    مگر اینکه نیروی خارجی به آن وارد شود.
    \label{thm:newton-first}
\end{theorem}

\mnote{قانون اول نیوتن را «قانون اینرسی» یا «قانون لختی» نیز می‌نامند.}

\subsection{قانون دوم}
\label{sec:physics-book-newton-second}

\begin{theorem}[قانون دوم نیوتن]
    شتاب جسم با نیروی خالص وارد بر آن
    متناسب و با جرم آن نسبت عکس دارد:
    \[
        \vec{F}_{\text{net}} = m \vec{a}
    \]
    \label{thm:newton-second}
\end{theorem}

طبق \cref{thm:newton-second}، اگر نیروی
خالص وارد بر جسمی صفر باشد، شتاب آن نیز
صفر است.

\begin{example}
    جسمی به جرم \lr{$2\,\text{kg}$} تحت
    نیروی \lr{$10\,\text{N}$} قرار دارد.
    شتاب آن برابر است با:
    \[
        a = \frac{F}{m} = \frac{10}{2} = 5\,\text{m/s}^2
    \]
    \label{ex:newton-second}
\end{example}

\subsection{قانون سوم}
\label{sec:physics-book-newton-third}

\begin{theorem}[قانون سوم نیوتن]
    برای هر کنش، واکنشی مساوی و در جهت
    مخالف وجود دارد:
    \[
        \vec{F}_{12} = -\vec{F}_{21}
    \]
    \label{thm:newton-third}
\end{theorem}

\mnote{قوانین نیوتن \cref{thm:newton-first,thm:newton-second,thm:newton-third} پایه‌ی مکانیک کلاسیک هستند.}

\section{معادلات حرکت}
\label{sec:physics-book-equations}

\subsection{حرکت با شتاب ثابت}
\label{sec:physics-book-constant-accel}

برای حرکت با شتاب ثابت \lr{$a$}، معادلات
حرکت به‌صورت زیر هستند:

\begin{align}
    v &= v_0 + at \label{eq:velocity} \\
    x &= x_0 + v_0 t + \frac{1}{2} a t^2 \label{eq:position} \\
    v^2 &= v_0^2 + 2a(x - x_0) \label{eq:velocity-position}
\end{align}

معادلات \cref{eq:velocity,eq:position,eq:velocity-position}
به‌عنوان معادلات سینماتیک شناخته می‌شوند.

\mnote{این معادلات فقط برای شتاب ثابت معتبرند. برای شتاب متغیر، باید از حساب دیفرانسیل استفاده کرد.}

\begin{example}
    خودرویی با سرعت اولیه‌ی
    \lr{$20\,\text{m/s}$} با شتاب
    \lr{$-5\,\text{m/s}^2$} ترمز می‌کند.
    مدت زمان تا توقف:
    \[
        t = \frac{v - v_0}{a} = \frac{0 - 20}{-5} = 4\,\text{s}
    \]
    \label{ex:braking}
\end{example}

\subsection{حرکت پرتابی}
\label{sec:physics-book-projectile}

برای حرکت پرتابی در دو بُعد، معادلات
به‌صورت زیر هستند:

\begin{align}
    x(t) &= v_0 \cos\theta \cdot t \\
    y(t) &= v_0 \sin\theta \cdot t - \frac{1}{2} g t^2
\end{align}

\mnote{در حرکت پرتابی، مؤلفه‌ی افقی سرعت ثابت است ولی مؤلفه‌ی عمودی تحت تأثیر گرانش تغییر می‌کند.}

\section{یکاها و ابعاد}
\label{sec:physics-book-units}

در فیزیک، استفاده‌ی درست از یکاها
ضروری است. جدول \cref{tab:si-units}
یکاهای اصلی \lr{SI} را نشان می‌دهد.

\begin{matintable}{یکاهای اصلی \lr{SI}}{tab:si-units}
\begin{tblr}{
    width = \textwidth,
    colspec = {l X[l] X[l] X[c]},
    hlines, vlines,
    colsep = 8pt,
    rowsep = 4pt,
    row{1} = {font=\bfseries},
}
کمیت & نام یکا & نماد & بُعد \\
طول & متر & \lr{m} & \lr{L} \\
جرم & کیلوگرم & \lr{kg} & \lr{M} \\
زمان & ثانیه & \lr{s} & \lr{T} \\
جریان & آمپر & \lr{A} & \lr{I} \\
دما & کلوین & \lr{K} & \lr{Θ} \\
مقدار ماده & مول & \lr{mol} & \lr{N} \\
شدت روشنایی & کندلا & \lr{cd} & \lr{J} \\
\end{tblr}
\end{matintable}

\mnote{کمیت‌های فرعی (مثل نیرو، انرژی، توان) از ترکیب یکاهای اصلی ساخته می‌شوند.}

\subsection{تحلیل ابعادی}
\label{sec:physics-book-dimensional}

\begin{example}
    بررسی ابعادی معادله‌ی \lr{$v = v_0 + at$}:
    \begin{itemize}
        \item \lr{$v$} و \lr{$v_0$}: \lr{$LT^{-1}$}.
        \item \lr{$at$}: \lr{$LT^{-2} \cdot T = LT^{-1}$}.
    \end{itemize}
    بنابراین دو طرف معادله، ابعاد یکسان
    دارند.
    \label{ex:dimensional-analysis}
\end{example}

\mnote{تحلیل ابعادی، ابزاری قدرتمند برای بررسی درستی معادلات فیزیکی است.}

\section{داده‌های تجربی}
\label{sec:physics-book-data}

جدول \cref{tab:free-fall-data} داده‌های
تجربی مربوط به سقوط آزاد را نشان می‌دهد.

\begin{matintable}{داده‌های تجربی سقوط آزاد}{tab:free-fall-data}
\begin{tblr}{
    width = \textwidth,
    colspec = {c c c},
    hlines, vlines,
    colsep = 8pt,
    rowsep = 4pt,
    row{1} = {font=\bfseries},
}
زمان \lr{(s)} & سرعت \lr{(m/s)} & ارتفاع \lr{(m)} \\
۰ & ۰٫۰ & ۰٫۰ \\
۰٫۵ & ۴٫۹ & ۱٫۲ \\
۱٫۰ & ۹٫۸ & ۴٫۹ \\
۱٫۵ & ۱۴٫۷ & ۱۱٫۰ \\
۲٫۰ & ۱۹٫۶ & ۱۹٫۶ \\
\end{tblr}
\end{matintable}

\mnote{داده‌های تجربی معمولاً با عدم قطعیت همراه هستند. در این جدول، مقادیر با یک رقم اعشار گزارش شده‌اند.}

\section{نمودارها}
\label{sec:physics-book-graphs}

نمودار \cref{fig:free-fall} داده‌های
سقوط آزاد را نمایش می‌دهد.

\begin{figure}[h]
\centering
\begin{latin}
\begin{tikzpicture}
  \begin{axis}[
    width=0.8\textwidth, height=6cm,
    xlabel={$t$ (s)},
    ylabel={$v$ (m/s)},
    grid=major,
    legend pos=north west
  ]
    \addplot[only marks, mark=*, color=matinblue, mark size=2.5pt]
        coordinates {(0,0) (0.5,4.9) (1,9.8) (1.5,14.7) (2,19.6)};
    \addlegendentry{|\pc{داده‌های تجربی}|}

    \addplot[matinred, thick, domain=0:2] {9.8*x};
    \addlegendentry{$v = gt$}
  \end{axis}
\end{tikzpicture}
\end{latin}
\caption{نمودار سرعت-زمان در سقوط آزاد}
\label{fig:free-fall}
\end{figure}

\mnote{نمودار \lr{\cref{fig:free-fall}} نشان می‌دهد که سرعت با زمان خطی افزایش می‌یابد و شیب خط، شتاب گرانش \lr{$g$} است.}

\section{تمرین‌ها}
\label{sec:physics-book-exercises}

\begin{exercise}
    جسمی از ارتفاع \lr{$45\,\text{m}$}
    رها می‌شود. با فرض \lr{$g = 10\,\text{m/s}^2$}،
    سرعت آن هنگام برخورد با زمین چقدر
    است؟
    \label{exc:free-fall}
\end{exercise}

\begin{solution}
    از معادله‌ی \cref{eq:velocity-position}:
    \[
        v^2 = v_0^2 + 2g h = 0 + 2 \cdot 10 \cdot 45 = 900
    \]
    بنابراین \lr{$v = 30\,\text{m/s}$}.
\end{solution}

\begin{exercise}
    خودرویی با سرعت \lr{$72\,\text{km/h}$}
    در حال حرکت است. راننده ترمز می‌کند و
    خودرو پس از \lr{$40\,\text{m}$} متوقف
    می‌شود. شتاب ترمز چقدر است؟
    \label{exc:braking-distance}
\end{exercise}

\begin{exercise}
    ابعاد کمیت \lr{$F = ma$} را با استفاده
    از یکاهای اصلی \lr{SI} بیان کنید.
    \label{exc:force-dimension}
\end{exercise}

\section{خلاصه}
\label{sec:physics-book-summary}

در این فصل، یک فصل کامل از یک کتاب
فیزیک را دیدیم که شامل:

\begin{itemize}
    \item \textbf{قضایا}: سه قانون نیوتن.
    \item \textbf{معادلات}: سینماتیک،
    حرکت پرتابی.
    \item \textbf{اثبات‌ها}: برای قضایا.
    \item \textbf{مثال‌ها}: با محاسبات عددی.
    \item \textbf{جداول}: یکاهای \lr{SI}،
    داده‌های تجربی.
    \item \textbf{نمودارها}: سرعت-زمان.
    \item \textbf{تمرین‌ها}: با راه‌حل.
\end{itemize}

\mnote{این فصل، الگویی برای نوشتن فصل‌های کتاب فیزیک شماست. می‌توانید ساختار آن را برای موضوعات دیگر (مثل الکترومغناطیس، ترمودینامیک) تطبیق دهید.}

این فصل، آخرین فصل از بخش پنجم (مثال‌های
کامل) بود. در بخش ششم (مرجع دستورات)،
مرجع کامل کلاس، محیط‌ها و دستورات ارائه
می‌شود.